a) Setiap matriks nilpoten
\mathl{\begin{pmatrix}  a   &   b   \\  c  & d \end{pmatrix}}{}
dapat dihubungkan dengan matriks nol melalui lintasan linear
\maabbeledisp {} {[0,1]} {N
} { t } {t \begin{pmatrix}  a   &   b   \\  c  & d \end{pmatrix}
} {,}
yang seluruhnya berada di dalam himpunan matriks nilpoten. Jadi $N$ terhubung lintasan dan, karena itu, terhubung.

b) Sebuah
$2 \times 2$-\definitionsverweis {matriks}{}{}
nilpoten tepat ketika jejak dan determinannya sama dengan $0$. Dengan demikian, himpunan $N$ semua matriks nilpoten dapat ditulis sebagai

\mathdisp {\left\{  (a,b,c,d)  \mid   a+d = 0 \text{ dan } ad-bc = 0         \right\}} { . }

Perhatikan pemetaan
\maabbeledisp {} {\R^4} {\R^2
} {(a,b,c,d)} { (a+d, ad-bc)
} {.}
Matriks Jacobi pemetaan ini ialah

\mathdisp {\begin{pmatrix}  1 &  0 &  0 &  1 \\  d  &  -c &  -b &  a  \end{pmatrix}} { . }

Pemetaan ini tidak
\definitionsverweis {reguler}{}{}
di titik nol, tetapi reguler di setiap titik lain pada serat $N$. Memang, jika
\mathl{P \in M}{}, maka dari syarat jejak dan determinan diperoleh

\mavergleichskettedisp
{\vergleichskette
{a^2
}
{ =} {-bc
}
{ } {
}
{ } {
}
{ } {
}
}
{}{}{}.
Karena itu,

\mavergleichskettedisp
{\vergleichskette
{b
}
{ =} { c
}
{ =} { 0
}
{ } {
}
{ } {
}
}
{}{}{}
langsung mengakibatkan

\mavergleichskettedisp
{\vergleichskette
{a
}
{ =} {b
}
{ =} { c
}
{ =} { d
}
{ =} { 0
}
}
{}{}{.}
Jadi matriks Jacobi mempunyai rank maksimal di semua titik $M$. Dengan

\mavergleichskettedisp
{\vergleichskette
{G
}
{ =} {\R^4 \setminus \{0\}
}
{ } {
}
{ } {
}
{ } {
}
}
{}{}{}
kita dapat memandang $M$ sebagai serat pemetaan yang dibatasi ke $G$, dan pemetaan itu reguler di seluruh serat. Oleh karena itu, berdasarkan
[[Satz über implizite Abbildungen/Faser ist Mannigfaltigkeit/Fakt|fakta tentang serat reguler]],
$M$ merupakan submanifold tertutup dari $G$.

Catatan edisi: Sumber menulis $a^2=bc$ setelah memakai $d=-a$. Tanda yang benar dari $ad-bc=0$ adalah $a^2=-bc$; argumen rank di atas menggunakan identitas yang telah diperbaiki ini.

c) Berdasarkan
[[Satz über implizite Abbildungen/Faser ist Mannigfaltigkeit/Fakt|fakta tentang serat reguler]],
dimensi $M$ ialah

\mavergleichskette
{\vergleichskette
{ 4-2
}
{ = }{ 2
}
{  }{
}
{    }{
}
{   }{
}
}
{}{}{.}

d) Tuliskan

\mavergleichskettedisp
{\vergleichskette
{M_+
}
{ =} { \left\{ (a,b,c,d) \in M  \mid   b>c        \right\}
}
{ } {
}
{ } {
}
{ } {
}
}
{}{}{}
dan

\mavergleichskettedisp
{\vergleichskette
{M_-
}
{ =} { \left\{ (a,b,c,d) \in M  \mid   b<c        \right\}
}
{ } {
}
{ } {
}
{ } {
}
}
{}{}{.}
Keduanya merupakan himpunan terbuka di $M$
\zusatzklammer {sebagai irisan $M$ dengan himpunan terbuka di $\R^4$ yang masing-masing ditentukan oleh $b>c$ dan $b<c$} {} {.}
Matriks
\mathkor {} {\begin{pmatrix}  0  &   1  \\  0 &  0 \end{pmatrix}} {dan} {\begin{pmatrix}  0  &   0  \\  1 &  0 \end{pmatrix}} {}
menunjukkan bahwa keduanya tidak kosong. Selain itu, keduanya menutupi seluruh $M$. Jika

\mavergleichskette
{\vergleichskette
{b
}
{ = }{ c
}
{  }{
}
{    }{
}
{   }{
}
}
{}{}{},
maka

\mavergleichskettedisp
{\vergleichskette
{ 0
}
{ =} { ad-bc
}
{ =} { -a^2 -b^2
}
{ } {
}
{ } {
}
}
{}{}{}
langsung memberikan

\mavergleichskette
{\vergleichskette
{a
}
{ = }{b
}
{ = }{ c
}
{ =   }{ d
}
{ =  }{ 0
}
}
{}{}{,}
dan titik ini tidak termasuk $M$. Jadi $M$ mempunyai tutupan oleh dua himpunan terbuka tak kosong yang saling lepas; akibatnya, $M$ tidak terhubung.

e) Untuk $M_+$, gunakan pemetaan
\maabbeledisp {\psi} {\R^2 \setminus \{(0,0)\}} { M_+
} {(a,u)} { (a, u + \sqrt{a^2+u^2}, u - \sqrt{a^2 +u^2},-a)
} {.}
Karena

\mavergleichskettealign
{\vergleichskettealign
{ad-bc
}
{ =} { -a^2 - \left(  u + \sqrt{a^2+u^2}  \right) \left( u - \sqrt{a^2 +u^2}  \right)
}
{ =} { -a^2 - \left(  u^2 - \left(  a^2+u^2 \right)   \right)
}
{ =} { 0
}
{ } {
}
}
{}
{}{}
syarat determinan terpenuhi, dan karena

\mavergleichskettedisp
{\vergleichskette
{b-c
}
{ =} { u + \sqrt{a^2+u^2} - \left(  u- \sqrt{a^2 +u^2})  \right)
}
{ =} { 2 \sqrt{a^2+u^2}
}
{ >} { 0
}
{ } {
}
}
{}{}{},
citranya terletak di $M_+$. Pemetaan
\maabbeledisp {\varphi} {M_+} { \R^2 \setminus \{(0,0)\}
} {(a,b,c,d)} { \left(  a  , \,   { \frac{  b+c }{  2 } }                   \right)
} {,}
merupakan invers dari pemetaan tersebut. Identitas

\mavergleichskettedisp
{\vergleichskette
{ \varphi \circ \psi
}
{ =} {
\operatorname{Id}_{ \R^2 \setminus \{(0,0)\}  }
}
{ } {
}
{ } {
}
{ } {
}
}
{}{}{}
jelas. Identitas lainnya mengikuti dari

\mavergleichskettealign
{\vergleichskettealign
{ { \frac{  b+c }{  2 } } + \sqrt{a^2 + \left( { \frac{  b+c }{  2 } }  \right)^2 }
}
{ =} { { \frac{  b+c }{  2 } } + { \frac{   1 }{  2 } }  \sqrt{4 a^2 + b^2 +c^2 +2bc }
}
{ =} { { \frac{  b+c }{  2 } } + { \frac{   1 }{  2 } }  \sqrt{ b^2 +c^2 -2bc }
}
{ =} { { \frac{  b+c }{  2 } } + { \frac{  b-c }{  2 } }
}
{ =} {b
}
}
{}
{}{}
dan

\mavergleichskettedisp
{\vergleichskette
{ { \frac{  b+c }{  2 } } - \sqrt{a^2 + \left(   { \frac{  b+c }{  2 } }   \right)^2 }
}
{ =} { { \frac{  b+c }{  2 } } - { \frac{  b-c }{  2 } }
}
{ =} { c
}
{ } {}
{ } {}
}
{}{}{.}
Kedua pemetaan kontinu, sehingga diperoleh suatu
\definitionsverweis {homeomorfisme}{}{.}
Untuk
\mathl{M_-}{}, tukarkan peran
\mathkor {} {b} {dan} {c} {.}

 [[Kategorie:Latexseite]]
